\documentclass[11pt]{amsart}

\usepackage[T1]{fontenc}
\usepackage{lmodern}
\usepackage{microtype}
\usepackage{amsmath,amssymb,amsthm}
\usepackage[hidelinks]{hyperref}

\newtheorem{theorem}{Theorem}
\newtheorem{proposition}[theorem]{Proposition}
\newtheorem{lemma}[theorem]{Lemma}
\newtheorem{corollary}[theorem]{Corollary}

\setlength{\emergencystretch}{2em}

\DeclareMathOperator{\CAB}{CAB}
\newcommand{\Dc}{D_{c}}

\title[Cyclic antibandwidth of $P_2\times P_2\times P_n$]{A Matching Labelling for the Cyclic
Antibandwidth of the \texorpdfstring{$2\times2\times n$}{2x2xn} Mesh}
\date{}
\subjclass[2020]{Primary 05C78; Secondary 05C76, 05C85, 68R10}
\keywords{cyclic antibandwidth, graph labelling, three-dimensional mesh, Cartesian product,
cylinder, extremal labelling, benchmark conjecture}

\hypersetup{
  pdftitle={A Matching Labelling for the Cyclic Antibandwidth of the 2x2xn Mesh}
}

\begin{document}

\begin{abstract}
We prove that the cyclic antibandwidth of the three-dimensional mesh $P_2\times P_2\times P_n$,
equivalently of the cylinder $C_4\,\square\,P_n$, is $2(n-1)$ for every integer $n\ge2$: one
closed-form labelling attains $2n-2$, and a three-line counting bound shows nothing attains more
for $n\ge3$, the endpoint $n=2$ being settled separately. In particular
$\CAB(P_{2\times2\times335}) = 668$ and $\CAB(P_{2\times2\times500}) = 998$; the identity is false
at $n=1$. This settles the first row of a table of conjectured three-dimensional mesh values
attributed to Bansal and Srivastava, which asserts
$\CAB(P_2\times P_2\times P_{n_3}) = 2(n_3-1)$ for $n_3 = 3,\dots,500$; two of those $498$ cells
already carried a certificate of optimality in the benchmark paper that reproduces the table, and
it is from that paper, not from the 2011 original, that we quote the row. The upper bound is also a
two-line consequence of the published two-dimensional mesh formula evaluated at $n_2 = 4$, which we
record; the matching labelling of the cylinder we do not claim as new, the two papers carrying that
formula not having been obtained in full.
\end{abstract}

\maketitle

\section{The statement}

Throughout, $G$ is a finite simple graph on $N = |V(G)|$ vertices. Following
\cite[\S2.1]{SAT2026}: a \emph{labelling} of $G$ is a bijection $f\colon V(G)\to\{1,\dots,N\}$; for
an edge $\{x,y\}$ the \emph{cyclic distance} is
\[
  \Dc(x,y) \;=\; \min\bigl(\,|f(x)-f(y)|,\; N-|f(x)-f(y)|\,\bigr);
\]
and
\[
  \CAB_f(G) \;=\; \min_{\{x,y\}\in E(G)} \Dc(x,y),
  \qquad
  \CAB(G) \;=\; \max_{f} \CAB_f(G).
\]
A \emph{three-dimensional mesh} $P_{n_1\times n_2\times n_3}$ is the Cartesian product
$P_{n_1}\times P_{n_2}\times P_{n_3}$ of three paths \cite[Appendix~A]{SAT2026}.

A table of conjectured exact $\CAB$ values for five families of three-dimensional meshes is
attributed to Bansal and Srivastava \cite{BansalSrivastava}. We did not read their 2011 paper; the
table is reproduced and attributed in \cite{SAT2026}, in the table
captioned ``Conjectured CAB value of 3D mesh instances'', and it is from there that we quote. In
the \LaTeX{} e-print source of
arXiv:2601.04239v1 its first data row reads
\begin{verbatim}
2 & 2 & $3$--$500$ & $2(n_3 - 1)$ \rule{0pt}{12pt} \\
\end{verbatim}
where \verb|\rule| is vertical spacing with no mathematical content. That is, the conjecture under
consideration is
\begin{equation}\label{eq:conj}
  \CAB(P_2\times P_2\times P_{n_3}) \;=\; 2(n_3-1)
  \qquad\text{for } n_3 = 3,4,\dots,500 .
\end{equation}
The authors of \cite{SAT2026} caution that a conjectured value ``may be less than the actual optimal
CAB value''. Theorem~\ref{thm:main} below proves the equality.

The benchmark of \cite{SAT2026} reports four instances of this family, $n_3 = 3$, $168$, $335$ and
$500$. At $n_3 = 3$ and $n_3 = 168$ its exact solvers certify the conjectured values $4$ and $334$
as optimal, so two of the $498$ cells of \eqref{eq:conj} carry a certificate of optimality there;
at $n_3 = 335$ and $n_3 = 500$ none of them certifies anything. For those two instances the values
$667$ and $997$ are the best entries recorded in the OPTSICOM cyclic antibandwidth repository
\cite{Optsicom}; they are that repository's entries, not a statement about every labelling anyone
has built.

\begin{theorem}\label{thm:main}
For every integer $n \ge 2$,
\[
  \CAB(P_2\times P_2\times P_n) \;=\; 2(n-1).
\]
In particular \eqref{eq:conj} holds at all $498$ of its tabulated cells, and for every $n_3 > 500$
as well. The identity is false at $n=1$, where $\CAB(P_2\times P_2\times P_1) = \CAB(C_4) = 1$.
\end{theorem}

\section{The graph, and a labelling that attains \texorpdfstring{$2n-2$}{2n-2}}

Write $G_n = P_2\times P_2\times P_n$ and index $V(G_n)$ by triples $(u,v,i)$ with $u,v\in\{0,1\}$
and $0\le i\le n-1$, two triples being adjacent exactly when they differ in one coordinate and
differ there by $1$. Thus
\[
  N \;=\; |V(G_n)| \;=\; 4n, \qquad |E(G_n)| \;=\; 8n-4,
\]
the edges falling into three families: $2n$ \emph{$v$-edges} $(u,v,i)(u,1-v,i)$, $2n$
\emph{$u$-edges} $(u,v,i)(1-u,v,i)$, and $4(n-1)$ \emph{$i$-edges} $(u,v,i)(u,v,i+1)$. Since
$P_2\times P_2$ is the $4$-cycle $(0,0)(1,0)(1,1)(0,1)$, the graph $G_n$ is the cylinder
$C_4\,\square\,P_n$. For $n\ge3$ its degree sequence consists of eight $3$'s and $4(n-2)$ $4$'s, so
$\Delta(G_n) = 4$.

\begin{proposition}\label{prop:lower}
For every integer $n\ge2$ the map
\begin{equation}\label{eq:f}
  f(u,v,i) \;=\; \bigl(1 + 2i + u\bigr) \;+\; 2n\cdot\bigl((u+v+i) \bmod 2\bigr)
\end{equation}
is a labelling of $G_n$, and $\CAB_f(G_n) = 2n-2$. Its edge distances are $2n$ on every $v$-edge,
$2n-1$ on every $u$-edge and $2n-2$ on every $i$-edge.
\end{proposition}

\begin{proof}
Put $g(u,i) = 1+2i+u$ and $b(u,v,i) = (u+v+i)\bmod2$, so that $f = g + 2n\,b$. The map
$(u,i)\mapsto g$ is a bijection $\{0,1\}\times\{0,\dots,n-1\}\to\{1,\dots,2n\}$, being the base-$2$
expansion, and for fixed $(u,i)$ the two choices of $v$ give the two values of $b$. Hence
$(u,v,i)\mapsto(g,b)$ is a bijection onto $\{1,\dots,2n\}\times\{0,1\}$ and
$(g,b)\mapsto g+2nb$ is a bijection onto $\{1,\dots,4n\}$. So $f$ is a labelling.

Every edge of $G_n$ changes exactly one coordinate by $1$, hence flips the parity $b$. Let
$d = |\Delta g|$ across an edge. Then $\Delta f = \Delta g \pm 2n$, so
$|\Delta f| \in \{2n-d,\,2n+d\}$, and since $d\le 2\le 2n$ both members lie in $[0,4n]$; as
$\min(t, 4n-t)$ takes the same value at $t = 2n-d$ and $t = 2n+d$, in either case
\[
  \Dc \;=\; 2n - d .
\]
Now $d = 0$ on a $v$-edge (the coordinate $v$ does not occur in $g$), $d = 1$ on a $u$-edge and
$d = 2$ on an $i$-edge, giving the three stated values. The three families have
$2n + 2n + 4(n-1) = 8n-4 = |E(G_n)|$ members, so they exhaust $E(G_n)$. For $n\ge2$ the $i$-edges
are present, so the minimum is $2n-2$.
\end{proof}

\section{The matching upper bound}

\begin{lemma}\label{lem:window}
Let $G$ be a graph on $N\ge2$ vertices and let $k$ be an integer with $1\le k\le N/2$. If
$\CAB(G)\ge k$ then every vertex of $G$ has degree at most $N-2k+1$.
\end{lemma}

\begin{proof}
Fix a labelling $f$ with $\CAB_f(G)\ge k$ and read labels in $\mathbb{Z}_N$. For $a\in\mathbb{Z}_N$
the residues at cyclic distance at most $k-1$ from $a$ are $a, a\pm1,\dots,a\pm(k-1)$, and these
$2k-1$ residues are distinct because $2(k-1) < N$. Hence exactly $N-(2k-1)$ residues lie at cyclic
distance $\ge k$ from $a$. Let $x$ be a vertex. Its $\deg(x)$ neighbours carry distinct labels, all
at cyclic distance $\ge k$ from $f(x)$, so $\deg(x)\le N-2k+1$.
\end{proof}

Equivalently $\CAB(G)\le\lfloor(N-\Delta(G)+1)/2\rfloor$ whenever the right-hand side is at most
$N/2$. This bound is elementary and is not claimed as new; see \S\ref{sec:print}.

\begin{theorem}\label{thm:upper}
$\CAB(G_n)\le 2n-2$ for every $n\ge3$.
\end{theorem}

\begin{proof}
Here $N = 4n$ and, since $n\ge3$, the vertex $(0,0,1)$ has the four distinct neighbours
$(1,0,1)$, $(0,1,1)$, $(0,0,0)$, $(0,0,2)$. Suppose $\CAB(G_n)\ge k$ for some $k\ge 2n-1$. As
$\CAB_f\ge k$ implies $\CAB_f\ge 2n-1$, we may take $k = 2n-1$, which satisfies
$k\le 2n = N/2$. Lemma~\ref{lem:window} then forces
\[
  4 \;\le\; N-2k+1 \;=\; 4n - 2(2n-1) + 1 \;=\; 3,
\]
a contradiction.
\end{proof}

Proposition~\ref{prop:lower} and Theorem~\ref{thm:upper} give
$\CAB(G_n) = 2(n-1)$ for all $n\ge3$, which is Theorem~\ref{thm:main} on the tabulated range and
beyond. The two endpoints are settled separately.

\begin{proposition}\label{prop:ends}
$\CAB(G_1) = \CAB(C_4) = 1$ and $\CAB(G_2) = \CAB(Q_3) = 2$.
\end{proposition}

\begin{proof}
For $n=1$, $N=4$ and the only pairs of residues at cyclic distance $2$ in $\mathbb{Z}_4$ are
$\{1,3\}$ and $\{2,4\}$. A labelling with $\CAB_f(C_4)\ge2$ would make each of the four edges of
$C_4$ such a pair, but the four edges are four distinct pairs and only two are available. Hence
$\CAB(C_4)=1$, while $2(n-1)=0$.

For $n=2$, $G_2 = Q_3$ is $3$-regular on $N=8$ vertices, and Lemma~\ref{lem:window} with $k=3$
gives no contradiction. Suppose $\CAB_f(Q_3)\ge3$. By the count in the proof of
Lemma~\ref{lem:window} exactly $8-2\cdot3+1 = 3$ residues lie at cyclic distance $\ge3$ from a
given $a\in\mathbb{Z}_8$, namely $a+3$, $a+4$, $a+5$. Every vertex has exactly three neighbours, so
those three residues are all used: $f$ is an isomorphism from $Q_3$ onto the circulant
$\mathrm{Circ}(8;\{\pm3,4\})$. Since $\gcd(3,8)=1$ the $\pm3$ edges of that circulant form the
single $8$-cycle $0,3,6,1,4,7,2,5$, and the four $4$-edges join the antipodal pairs of this cycle,
so the circulant is the M\"obius ladder on eight vertices. It contains the $5$-cycle
$0-3-6-1-4-0$, with successive differences $3,3,3,3,4$, and is therefore not bipartite, whereas
$Q_3$ is. This contradiction gives $\CAB(Q_3)\le2$, and Proposition~\ref{prop:lower} gives
$\CAB(Q_3)\ge2$.
\end{proof}

\section{Two instances written out}

For $n=3$ ($N=12$, $20$ edges) the labelling \eqref{eq:f} is, as \texttt{label = (u,v,i)},
\begin{verbatim}
 1=(0,0,0)   2=(1,1,0)   3=(0,1,1)   4=(1,0,1)
 5=(0,0,2)   6=(1,1,2)   7=(0,1,0)   8=(1,0,0)
 9=(0,0,1)  10=(1,1,1)  11=(0,1,2)  12=(1,0,2)
\end{verbatim}
with edge distance multiset $\{4^{8},\,5^{6},\,6^{6}\}$ and minimum $4$. This agrees with the
independently certified optimum $4^{*}$ of \cite{SAT2026}. For $n=4$ ($N=16$, $28$ edges) it is
\begin{verbatim}
 1=(0,0,0)   2=(1,1,0)   3=(0,1,1)   4=(1,0,1)
 5=(0,0,2)   6=(1,1,2)   7=(0,1,3)   8=(1,0,3)
 9=(0,1,0)  10=(1,0,0)  11=(0,0,1)  12=(1,1,1)
13=(0,1,2)  14=(1,0,2)  15=(0,0,3)  16=(1,1,3)
\end{verbatim}
with edge distance multiset $\{6^{12},\,7^{8},\,8^{8}\}$ and minimum $6 = 2(4-1)$.

At the two largest cells of \eqref{eq:conj} the same formula gives
$\CAB(P_{2\times2\times335}) = 668$ on $1340$ vertices and $\CAB(P_{2\times2\times500}) = 998$ on
$2000$ vertices, one above each of the entries $667$ and $997$ recorded in \cite{Optsicom}.

\section{What is already in print, and what is not settled}\label{sec:print}

\subsection*{The upper bound is a corollary of a published two-dimensional formula}
The two-dimensional mesh optimum $\CAB(P_{n_1}\times P_{n_2}) = n_2(n_1-1)/2$ is attributed to
\cite{STV2005} and \cite{RSSTV2009}; the OPTSICOM repository \cite{Optsicom} records it as an exact
equality in every case except $n_1$ even with $n_2$ odd, where only the bracket
$\lfloor n_2(n_1-1)/2\rfloor\le\CAB\le\lceil n_2(n_1-1)/2\rceil$ is stated as proved. At $n_2 = 4$
the exceptional case does not arise and the value is the integer $2(n_1-1)$, so
$\CAB(P_n\times P_4) = 2(n-1)$ is in print.
The grid $P_n\times P_4$ is a spanning subgraph of $C_4\,\square\,P_n = G_n$ --- the cylinder has
exactly $n$ more edges, one per layer --- and $\CAB$ cannot increase when edges are added on a fixed
vertex set. Hence Theorem~\ref{thm:upper} follows from that formula for $n\ge4$. We give
Lemma~\ref{lem:window} because it is self-contained, because we could not read \cite{STV2005} or
\cite{RSSTV2009} in full (see below), and because it also covers $n=3$. What is left beyond
that formula is Proposition~\ref{prop:lower}, the matching labelling of the \emph{cylinder}, which
we do not claim as new either --- whether it is already implicit in \cite{STV2005} or
\cite{RSSTV2009} could not be checked. Stated at its
sharpest, what is proved here is that closing the four-vertex direction of the $4\times n$ grid
into a $4$-cycle does not decrease its cyclic antibandwidth.

\subsection*{Lemma~\ref{lem:window} is elementary, and is not tight}
It is not claimed as new, and it is not generally tight: for
$P_3\times P_3\times P_3$ it gives $\lfloor(27-6+1)/2\rfloor = 11$, while \cite{SAT2026} certifies
the optimum $9$. So nothing here transfers to the other four rows of the conjectured table.

\subsection*{Adjacent work, and what it does not settle}
The nearest published items are \cite{STV2005}, which introduces the cyclic antibandwidth problem,
and \cite{RSSTV2009} on antibandwidth and cyclic antibandwidth of meshes and hypercubes. Of the
latter we could read the abstract but not the text; it states that the authors ``solve the
antibandwidth problem precisely for two-dimensional meshes, tori and estimate the antibandwidth
value for hypercubes up to the third-order term'', and that for the cyclic antibandwidth they
``prove basic facts and exact results for the same product graphs as for the antibandwidth''. The
exact cyclic values there are therefore for two-dimensional meshes and tori; $C_4\,\square\,P_n$ is
neither, and we did not find the value proved here recorded in any source we could reach. That is
not a novelty claim, and the unread texts bound it in two specific ways. First, the published proof of
$\CAB(P_n\times P_4) = 2(n-1)$ must exhibit an extremal labelling of the $4\times n$ grid; if that
labelling happens also to respect the $n$ wrap edges of the cylinder, then
Proposition~\ref{prop:lower} is an observation about published work rather than a new construction.
Second, since $C_4\,\square\,C_n$ is $G_n$ together with four further edges, and adding edges cannot
increase $\CAB$, a published toroidal value $2n-2$ at $C_4\,\square\,C_n$ would give
Proposition~\ref{prop:lower} at that $n$. Neither possibility has been checked either way, and
until \cite{RSSTV2009} is read the lower bound should not be described as new.
\cite{TorokVrto2010}, on the antibandwidth of three-dimensional meshes, was also not read.

\subsection*{Sources not read, and what rests on them}
\cite{BansalSrivastava} is not openly available and was not read. The conjecture is quoted from the
e-print of \cite{SAT2026}, which reproduces and attributes the table. If the 2011 wording differs
--- a different range, or a construction already attached --- that is not accounted for here. The
version of record of \cite{SAT2026} (\emph{Comput. Optim. Appl.}, published 21 August 2026) has no
openly readable location either; every quotation above is from arXiv:2601.04239v1 and could differ
there.

\subsection*{Scope}
The tabulated row is $n_3\in[3,500]$, and Theorem~\ref{thm:main} covers it together with every
larger $n_3$ and the endpoint $n=2$; it fails at $n=1$. The two cells that already carried
certificates of optimality in \cite{SAT2026} are used above as external checks; nothing above bears
on the other four rows of the conjectured table.
The accompanying program re-checks the printed labellings and, cell by cell over
$n = 2,\dots,500$, the bijectivity, degree multiset and edge-distance multiset on which
Proposition~\ref{prop:lower} and Lemma~\ref{lem:window} are used; beyond $n = 500$ the statement of
Theorem~\ref{thm:main} rests on those two proofs alone. The program does not fetch any of the
external documents quoted here and cannot bear on the priority question above.

\begin{thebibliography}{9}

\bibitem{BansalSrivastava}
R.~Bansal and K.~Srivastava,
\emph{A memetic algorithm for the cyclic antibandwidth maximization problem},
Soft Comput. \textbf{15} (2011), no.~2, 397--412.
\href{https://doi.org/10.1007/s00500-009-0538-6}{doi:10.1007/s00500-009-0538-6}

\bibitem{Optsicom}
OPTSICOM project,
\emph{Cyclic antibandwidth problem} (instances, best-known values and known optima),
\href{https://grafo.etsii.urjc.es/optsicom/cab/}{grafo.etsii.urjc.es/optsicom/cab/}.

\bibitem{RSSTV2009}
A.~Raspaud, H.~Schr\"oder, O.~S\'ykora, L.~T\"or\"ok and I.~Vr\v{t}o,
\emph{Antibandwidth and cyclic antibandwidth of meshes and hypercubes},
Discrete Math. \textbf{309} (2009), no.~11, 3541--3552.
\href{https://doi.org/10.1016/j.disc.2007.12.058}{doi:10.1016/j.disc.2007.12.058}

\bibitem{STV2005}
O.~S\'ykora, L.~T\"or\"ok and I.~Vr\v{t}o,
\emph{The cyclic antibandwidth problem},
Electron. Notes Discrete Math. \textbf{22} (2005), 223--227.
\href{https://doi.org/10.1016/j.endm.2005.06.030}{doi:10.1016/j.endm.2005.06.030}

\bibitem{SAT2026}
H.~Truong Xuan and K.~To Van,
\emph{Solving cyclic antibandwidth problem by SAT},
\href{https://arxiv.org/abs/2601.04239v1}{arXiv:2601.04239v1}, 2026;
Comput. Optim. Appl., published 21 August 2026,
\href{https://doi.org/10.1007/s10589-026-00819-8}{doi:10.1007/s10589-026-00819-8}

\bibitem{TorokVrto2010}
L.~T\"or\"ok and I.~Vr\v{t}o,
\emph{Antibandwidth of three-dimensional meshes},
Discrete Math. \textbf{310} (2010), no.~3, 505--510.
\href{https://doi.org/10.1016/j.disc.2009.03.029}{doi:10.1016/j.disc.2009.03.029}

\end{thebibliography}

\end{document}
